\section{Study Comparison Following Webster and Watson}
Following the concept matrix approach proposed by \parencite{websterAnalyzingPrepareFuture2002}, this section presents a structured comparison of the 13 identified studies on wearable seizure detection and forecasting. The matrix organizes studies along key dimensions relevant to clinical translation and technical implementation.

\subsection{Concept Matrix: Study Overview}

Table~\ref{tab:study-matrix} provides a comprehensive comparison of all included studies, organized by primary objective, methodology, and outcomes. The matrix enables cross-study analysis of trends in sensor modalities, algorithmic approaches, and validation strategies across the 2013--2026 period.

\begin{landscape}
\begin{table}[htbp]
\centering
\small
\caption{Webster \& Watson Concept Matrix: Wearable Seizure Detection and Forecasting Studies (2013--2026)}
\label{tab:study-matrix}
\setlength{\tabcolsep}{4pt}
\renewcommand{\arraystretch}{1.15}
\resizebox{\linewidth}{!}{%
\begin{tabular}{@{}llllllllllp{3cm}@{}}
\toprule
Study & Objective & Design & Sample & Device & Loc. & Algorithm & Sens & Spec & FAR & Key Limitation \\
\midrule

\textbf{Detection Studies} \\
\addlinespace

\textbf{Spahr et al. 2025} & Convulsive detection multi-center & Prospective & n=384, 49 CSs & Empatica E4 (ACC) & Wrist & Ensemble 1D CNN & 96\% & -- & $<$1/8d & EMU, offline \\
\addlinespace

\textbf{Reintjes et al. 2025} & ECG anomaly detection & Retro & n=120, 856 szs & Single-lead ECG & Chest & Matrix Profile & 98.16\% & -- & 1.91/h & High FAR \\
\addlinespace

\textbf{Fine 2025} & Tonic detection & Phase 1 & n=15/3 & 6-axis band & Wrist & ANN (594 feat) & 100\% & -- & 0.23/n; 0.023/h & Small test \\
\addlinespace

\textbf{Dong et al. 2026} & Nocturnal severe szs & Prospective & n=68, 1846 szs & NightWatch & Arm & CNN-LSTM & 71.6\% & -- & 0.165/h & Single-center \\
\addlinespace

\textbf{Wang et al. 2025} & Sz detection (acc: 97.8\%) & Retro & n=28, 62 szs & Biovital-P1 & Wrist & LSTM+SVM/LDA & -- & -- & 8.5/24h & Hospital only \\
\addlinespace

\textbf{Singh Rathore 2024} & Multimodal & Single pt & 1 pt, 6 hrs & EDA+ACC+HR & -- & MLP, SVM, KNN & 96.8\% & -- & -- & Single patient \\
\addlinespace

\textbf{Elemam et al. 2025} & Q+HRV+Audio & Cross-sect & Q:198, HRV:30 & Camera PPG & Finger & CNN, rule based & 95.1\% & 97.1\% & -- & Small samples \\
\addlinespace

\textbf{Borujeny 2013} & Wireless ACC & 3-pt & 3 pts, 20 szs & MICAz ACC & Arm/thigh & KNN k=5 & 100\% & -- & 0 & Very small \\
\addlinespace

\textbf{Forecasting Studies} \\
\addlinespace

\textbf{Vieluf et al. 2025} & Diary+wearable & Retro & n=70, 5437 d & Embrace & Wrist & DNN+harmonic & 82\% & 67\% & -- & Daily res only \\
\addlinespace

\textbf{Meisel et al. 2020} & LOSO forecast & LOSO CV & n=69, 452 szs & Empatica E4 & Wrist & LSTM+1DConv & 51.2\% & -- & TiW 43.7\% & Pediatric \\
\addlinespace

\textbf{Stirling 2021} & HR cycle forecast & Retro+pseudo & n=11, 136 szs & Fitbit & Wrist & LSTM+RF+LR & -- & -- & AUC 0.74 & Self-report \\
\addlinespace

\textbf{Nasseri 2021} & Ambulatory forecast & Retro & n=6 & Empatica E4 & Wrist & LSTM 4-layer & AUC 0.75 & -- & TiW 0.9--7.2h/d & RNS required \\
\addlinespace

\textbf{Ode et al. 2023} & RRI prediction & Retro & n=66, 85 szs & ECG & -- & Self-Att AE & 74\% & -- & 0.85/h & FPs in some pts \\
\bottomrule
\end{tabular}%
}
\par\medskip
\footnotesize
\textbf{Note:} ACC = accelerometer; ANN = artificial neural network; AE = autoencoder; CNN = convolutional neural network; CS = convulsive seizure; DNN = deep neural network; ECG = electrocardiogram; EDA = electrodermal activity; HR = heart rate; HRV = HR variability; LOSO = leave-one-subject-out; LSTM = long short-term memory; PPG = photoplethysmography; RRI = R-R interval; Sens = sensitivity; Spec = specificity; sz(s) = seizure(s); FPR = false positive rate; FA = false alarm; TiW = time in warning; -- = not reported.
\end{table}
\end{landscape}

\subsection{Dimensional Analysis}

\subsubsection{Sensor Modality Trends}
\begin{itemize}
  \item \textbf{Wrist-worn accelerometry} dominates (8/13 studies), leveraging commercial devices (Empatica E4, Fitbit, Embrace).
  \item \textbf{Cardiac signals} (ECG, PPG, HRV) show increasing adoption for preictal detection, with 6 studies incorporating autonomic measures.
  \item \textbf{Multi-modal fusion} (ACC + EDA + PPG) appears in 4 recent studies, reflecting recognition that single modalities have limited specificity.
\end{itemize}

\subsubsection{Algorithmic Evolution}
\begin{itemize}
  \item \textbf{2013:} Traditional ML (KNN, ANN) with handcrafted features \parencite{borujenyDetectionEpilepticSeizure2013}.
  \item \textbf2020--2021: Emergence of LSTM and ensemble methods \parencite{meiselMachineLearningWristband2020,stirlingForecastingSeizureLikelihood2021}.
  \item \textbf{2023--2026:} Deep learning dominance with CNNs, attention mechanisms, and anomaly detection approaches \parencite{reintjesECGBasedDetectionEpileptic2025,odeDevelopmentEpilepticSeizure2023,spahrDeepLearningbasedDetection2025}.
\end{itemize}

\subsubsection{Validation Rigor}
\begin{itemize}
  \item \textbf{Patient-level preservation:} Only 4 studies use LOSO or patient-independent splits.
  \item \textbf{Prospective validation:} Only 3 studies \parencite{spahrDeepLearningbasedDetection2025,dongTwoLayerEnsembleMethod2022,elemamAutomatedValidatedTool2025} include prospective phases.
  \item \textbf{Real-world testing:} Limited to 2 studies \parencite{vielufSeizureMonitoringCombined2025,nasseriAmbulatorySeizureForecasting2021}; most remain in EMU settings.
\end{itemize}

\subsubsection{Performance Metrics Reporting}
Table~\ref{tab:metrics-summary} summarizes the heterogeneity in performance metric reporting, complicating cross-study comparison. Sensitivity ranges from 51.2\% (forecasting mean) to 100\% (tonic detection), while FPR reporting varies from per-hour to per-night, with several studies omitting this critical metric entirely.

\begin{table}[htbp]
\centering
\small
\caption{Performance Metrics Reporting Across Studies}
\label{tab:metrics-summary}
\resizebox{\textwidth}{!}{
\begin{tabular}{lccc}
\toprule
\textbf{Metric} & \textbf{Detection (n=9)} & \textbf{Forecasting (n=4)} & \textbf{Reported By} \\
\midrule
Sensitivity/Recall & 8/9 & 2/4 & Reintjes, Fine, Spahr, Dong, Wang, \\
& & & Singh Rathore, Elemam, Borujeny, Meisel, Ode \\
Specificity & 1/9 & 0/4 & Elemam only \\
FPR/FA rate & 7/9 & 2/4 & Spahr, Reintjes, Fine, Dong, Wang, \\
& & & Ode, Meisel, Nasseri, Stirling \\
AUC-ROC & 1/9 & 2/4 & Reintjes, Nasseri, Stirling \\
Accuracy & 2/9 & 0/4 & Wang, Singh Rathore \\
Precision & 0/9 & 1/4 & Ode (reported as precision) \\
\bottomrule
\end{tabular}
}
\end{table}

\subsection{Synthesis and Gaps}

The concept matrix reveals several patterns:

\begin{enumerate}
  \item \textbf{Clinical translation gap:} Despite high sensitivities ($>$90\% in several detection studies), FPR remains clinically unacceptable for standalone deployment.
  \item \textbf{Personalization vs. generalization:} LOSO approaches achieve lower but more realistic performance (51.2\%) compared to patient-specific tuning, raising questions about generalizability.
  \item \textbf{Diary integration:} \parencite{vielufSeizureMonitoringCombined2025} demonstrates that hybrid approaches (wearable + clinical data) may outperform wearables alone, suggesting a direction for ambulatory systems.
  \item \textbf{Metric heterogeneity:} Inconsistent reporting (sensitivity vs. accuracy vs. AUC) impedes meta-analysis and clinical benchmarking.
\end{enumerate}
